\documentclass[10pt,a4paper]{amsart}
\usepackage{amsmath,amssymb,amsthm,graphicx}
\usepackage{verbatim}
\newenvironment{reference}{\comment}{\endcomment}
%\newenvironment{reference}{}{}
\newenvironment{slogan}{\comment}{\endcomment}
\newenvironment{history}{\comment}{\endcomment}
% For commutative diagrams we use Xy-pic
\usepackage[all]{xy}

% We use 2cell for 2-commutative diagrams.
\xyoption{2cell}
% Theorem environments.
%
\theoremstyle{plain}
\newtheorem{theorem}[subsection]{Theorem}
\newtheorem{proposition}[subsection]{Proposition}
\newtheorem{lemma}[subsection]{Lemma}

\theoremstyle{definition}
\newtheorem{definition}[subsection]{Definition}
\newtheorem{example}[subsection]{Example}
\newtheorem{exercise}[subsection]{Exercise}
\newtheorem{situation}[subsection]{Situation}

\theoremstyle{remark}
\newtheorem{remark}[subsection]{Remark}
\newtheorem{remarks}[subsection]{Remarks}

\numberwithin{equation}{subsection}

% Macros
%
\def\lim{\mathop{\mathrm{lim}}\nolimits}
\def\colim{\mathop{\mathrm{colim}}\nolimits}
\def\Spec{\mathop{\mathrm{Spec}}}
\def\Hom{\mathop{\mathrm{Hom}}\nolimits}
\def\Ext{\mathop{\mathrm{Ext}}\nolimits}
\def\SheafHom{\mathop{\mathcal{H}\!\mathit{om}}\nolimits}
\def\SheafExt{\mathop{\mathcal{E}\!\mathit{xt}}\nolimits}
\def\Sch{\mathit{Sch}}
\def\Mor{\mathop{\mathrm{Mor}}\nolimits}
\def\Ob{\mathop{\mathrm{Ob}}\nolimits}
\def\Sh{\mathop{\mathit{Sh}}\nolimits}
\def\NL{\mathop{N\!L}\nolimits}
\def\CH{\mathop{\mathrm{CH}}\nolimits}
\def\proetale{{pro\text{-}\acute{e}tale}}
\def\etale{{\acute{e}tale}}
\def\QCoh{\mathit{QCoh}}
\def\Ker{\mathop{\mathrm{Ker}}}
\def\Im{\mathop{\mathrm{Im}}}
\def\Coker{\mathop{\mathrm{Coker}}}
\def\Coim{\mathop{\mathrm{Coim}}}

% Boxtimes
%
\DeclareMathSymbol{\boxtimes}{\mathbin}{AMSa}{"02}

%
% Macros for moduli stacks/spaces
%
\def\QCohstack{\mathcal{QC}\!\mathit{oh}}
\def\Cohstack{\mathcal{C}\!\mathit{oh}}
\def\Spacesstack{\mathcal{S}\!\mathit{paces}}
\def\Quotfunctor{\mathrm{Quot}}
\def\Hilbfunctor{\mathrm{Hilb}}
\def\Curvesstack{\mathcal{C}\!\mathit{urves}}
\def\Polarizedstack{\mathcal{P}\!\mathit{olarized}}
\def\Complexesstack{\mathcal{C}\!\mathit{omplexes}}
% \Pic is the operator that assigns to X its picard group, usage \Pic(X)
% \Picardstack_{X/B} denotes the Picard stack of X over B
% \Picardfunctor_{X/B} denotes the Picard functor of X over B
\def\Pic{\mathop{\mathrm{Pic}}\nolimits}
\def\Picardstack{\mathcal{P}\!\mathit{ic}}
\def\Picardfunctor{\mathrm{Pic}}
\def\Deformationcategory{\mathcal{D}\!\mathit{ef}}
\usepackage{cleveref}
\crefname{originallemma}{original lemma}{original lemmas}
\Crefname{originallemma}{Original lemma}{Original lemmas}
\makeatletter
\newlabel{lemma-finite-type-surjective-on-stalk}{{lemma-finite-type-surjective-on-stalk}{source}{}{sourcepointer-1}{}}
\newlabel{lemma-finite-type-surjective-on-stalk@cref}{{[originallemma][0][]lemma-finite-type-surjective-on-stalk}{[1][0][]source}{}{}{}}
\makeatother
\crefname{originallemma}{original lemma}{original lemmas}
\Crefname{originallemma}{Original lemma}{Original lemmas}
\makeatletter
\newlabel{lemma-kernel-surjection-finite-free-onto-finite-presentation}{{lemma-kernel-surjection-finite-free-onto-finite-presentation}{source}{}{sourcepointer-2}{}}
\newlabel{lemma-kernel-surjection-finite-free-onto-finite-presentation@cref}{{[originallemma][0][]lemma-kernel-surjection-finite-free-onto-finite-presentation}{[1][0][]source}{}{}{}}
\makeatother
\crefname{originallemma}{original lemma}{original lemmas}
\Crefname{originallemma}{Original lemma}{Original lemmas}
\makeatletter
\newlabel{lemma-finite-type-stalk-zero}{{lemma-finite-type-stalk-zero}{source}{}{sourcepointer-3}{}}
\newlabel{lemma-finite-type-stalk-zero@cref}{{[originallemma][0][]lemma-finite-type-stalk-zero}{[1][0][]source}{}{}{}}
\makeatother
\title{A finitely presented module whose stalk at a point is free of rank r is free of that rank on a neighbourhood}
\author{The Stacks Project Authors}
\date{stacks-project:modules.tex:local-snapshot-20261012; GNU FDL 1.2 or later}
\begin{document}
\maketitle
\noindent\textit{Editorial note.} All delivered mathematical text is the author's own, copied from the pinned Stacks Project chapter source sources/stacks-modules/pinned-original-source.tex; no statement or proof step was rewritten, reordered or removed.\par\medskip
\noindent\textit{Editorial note.} Preamble provenance: the Stacks Project keeps one single shared preamble for every chapter; this card takes it directly from the staged sources/stacks-modules/preamble.tex as three verbatim spans (lines 8-12, 14-18 and 154-217), so no derived file was written and nothing was added or redefined. The wrapper class and the base packages are supplied by the assembler.\par\medskip
\noindent\textit{Editorial note.} Bibliography provenance: neither the counted statement nor its proof contains a citation (this lane rescans the whole span before writing the card), so the card carries no bibliography and no entry was reconstructed.\par\medskip
\noindent\textit{Editorial note.} Selection provenance: this card was opened under the file-based control contract (CONTROL.md), with the mechanical pre-card gate run on its statement and proof interval before the card was written.\par\medskip
\noindent\textit{Allowed standard prerequisites (editorial).}\begin{itemize}
\item The lifting of finitely many elements of a stalk to sections over a neighbourhood, and the spreading of surjectivity on a stalk to a neighbourhood. Used inside the counted proof (source lines 1763-1768), the second by a pointer to a lemma of this chapter. Produces a map from a finite free module onto the given module over a neighbourhood.
\item The finiteness of the kernel of a surjection from a finite free module onto a finitely presented module, and the criterion for a module of finite type whose stalk at a point vanishes to vanish on a neighbourhood. Both are invoked inside the counted proof (source lines 1769-1772) as pointers to lemmas of this chapter. Shows that the kernel of the constructed map vanishes after shrinking, whence the map is an isomorphism.
\end{itemize}\medskip
\ifdefined\hypertarget\hypertarget{sourcepointer-1}{}\fi
\noindent\textit{Source reference.} Original lemma lemma-finite-type-surjective-on-stalk: sources/stacks-modules/pinned-original-source.tex, source lines 1125--1125. Navigation pointer only: the counted text refers to the Stacks item lemma-finite-type-surjective-on-stalk, which belongs to another chapter of the project and is not part of this card. The pointer carries the target's own label as its number because this chapter's staged snapshot does not contain the other chapter's tag table; its content is not reproduced here.\par\medskip
\ifdefined\hypertarget\hypertarget{sourcepointer-2}{}\fi
\noindent\textit{Source reference.} Original lemma lemma-kernel-surjection-finite-free-onto-finite-presentation: sources/stacks-modules/pinned-original-source.tex, source lines 1769--1769. Navigation pointer only: the counted text refers to the Stacks item lemma-kernel-surjection-finite-free-onto-finite-presentation, which belongs to another chapter of the project and is not part of this card. The pointer carries the target's own label as its number because this chapter's staged snapshot does not contain the other chapter's tag table; its content is not reproduced here.\par\medskip
\ifdefined\hypertarget\hypertarget{sourcepointer-3}{}\fi
\noindent\textit{Source reference.} Original lemma lemma-finite-type-stalk-zero: sources/stacks-modules/pinned-original-source.tex, source lines 1140--1140. Navigation pointer only: the counted text refers to the Stacks item lemma-finite-type-stalk-zero, which belongs to another chapter of the project and is not part of this card. The pointer carries the target's own label as its number because this chapter's staged snapshot does not contain the other chapter's tag table; its content is not reproduced here.\par\medskip
\begin{lemma}
\label{lemma-finite-presentation-stalk-free}
Let $(X, \mathcal{O}_X)$ be a ringed space. Let $\mathcal{F}$ be
a finitely presented $\mathcal{O}_X$-module. Let $x \in X$ such that
$\mathcal{F}_x \cong \mathcal{O}_{X, x}^{\oplus r}$. Then there exists
an open neighbourhood $U$ of $x$ such that
$\mathcal{F}|_U \cong \mathcal{O}_U^{\oplus r}$.
\end{lemma}

\begin{proof}
Choose $s_1, \ldots, s_r \in \mathcal{F}_x$ mapping to a basis of
$\mathcal{O}_{X, x}^{\oplus r}$ by the isomorphism. Choose an open
neighbourhood $U$ of $x$ such that $s_i$ lifts to $s_i \in \mathcal{F}(U)$.
After shrinking $U$ we see that the induced map
$\psi : \mathcal{O}_U^{\oplus r} \to \mathcal{F}|_U$ is surjective
(Lemma \ref{lemma-finite-type-surjective-on-stalk}).
By Lemma \ref{lemma-kernel-surjection-finite-free-onto-finite-presentation}
we see that $\Ker(\psi)$ is of finite type.
Then $\Ker(\psi)_x = 0$ implies that $\Ker(\psi)$ becomes zero
after shrinking $U$ once more (Lemma \ref{lemma-finite-type-stalk-zero}).
\end{proof}

\end{document}
